\documentclass[a4paper]{article}
\input{../../../../newpreamble.tex}
\title{220: Midterm 1 Study Guide}
\author{Lance Remigio}
\begin{document}
\maketitle

\section*{Facts}

\subsection*{T1: Topological Spaces}

\begin{eg}
    \begin{itemize}
    \item If a topological space \( \mathcal{T} \) is \textbf{metrizable}, then the metric that induces \( \mathcal{T} \) is not unique.
    \item The coarsest topology is the \textbf{indiscrete topology} and the finest topology is the power set of a given nonempty set.
    \item For any topology \( \mathcal{T} \), we have the following inclusions:
        \[   \{ \emptyset, X  \}  \subseteq \mathcal{T} \subseteq  \mathcal{P}(x)  \]
    \item \( \mathcal{T} = \{ (a,b) : a,b \in \R , a < b  \}  \cup \{ \emptyset, \R  \}  \) is not a topology on \( \R  \) because \( (0,1)  \) and \( (2,3) \) are in \( \mathcal{T} \) but their union \( (0,1) \cup (2,3) \) is not in the form \( (a,b) \) and so \( (0,1) \cup (2,3) \notin \mathcal{T} \).
    \end{itemize}
\end{eg}

\subsection*{T2: Bases For Topological Spaces}

\begin{eg}
\begin{itemize}
    \item If \( (X,\mathcal{T}) \) is a topological space, then \( \mathcal{T} \) is basis for itself.
    \item In a metric space, the collection of all open balls is a basis for the topology induced by the metric.
    \item If \( X  \) is nonempty, the collection of all singletons is a basis for the discrete topology on \( X  \).
    \item The open balls generated by \( {d}_{2} \) and \( {d}_{\infty } \) metric form a basis for the Euclidean topology on \( \R^{2} \).
    \item Consider the following bases
        \[  {\mathcal{B}}_{E} = \{ (a,b) : a < b  \}, \ \ {\mathcal{B}}_{L} = \{ [a,b) : a < b  \}, \ \ {\mathcal{B}}_{K} = \mathcal{B} \cup \{ (a,b) \setminus  K : a < b \}     \]
        where \( K = \{ 1/n : n \in \N  \}  \). These bases generate the Euclidean topology \( {\mathcal{T}}_{E} \), the Lower Limit Topology \( {\mathcal{T}}_{L} \), and the \( K \)-topology \( {\mathcal{T}}_{K} \).

        The lower limit topology and the \( K \)-topology are strictly finer than the Euclidean one, but are not comparable to each other.
\end{itemize}    
\end{eg}

\begin{remark}
   Let \( X  \) be a nonempty set. If a collection of subsets \( \mathcal{B} \) of \( X \) satisfies the following properties:
   \begin{enumerate}
       \item[(i)] \( \bigcup_{ B \in \mathcal{B} }^{  } B = X  \)
        \item[(ii)] For all \( A,B \in \mathcal{B} \), then \( A \cap B  \) can be written as a union of elements in \( \mathcal{B} \)
   \end{enumerate}
   Then there exists a unique topology \( \mathcal{T} \) such that \( \mathcal{B} \) generates \( \mathcal{T} \).
\end{remark}

\begin{remark}
    In general, we can determine whether a collection of subsets of \( X  \) forms a basis for the topology. However, if the collection happens to satisfies the above properties, then it is a basis for some topology that is unique.
\end{remark}

\subsection*{T3: Neighborhoods, Interior}

\begin{eg}[Neighborhoods That Are Open/Not Open]
    \begin{itemize}
        \item For all \( r > 0  \), \( (x-r, x + r) \) is an open neighborhood containing \( x  \) since it is an open set containing \( x  \). On the other hand, \( [x-r, x+r) \) is not open and contains the open set \( (x-r, x + r) \). Thus, \( [x-r,x+r) \) is a neighborhood but is not open. 
        \item For all \( r > 0  \), \( [x,x + r) \) is not a neighborhood of \( x  \), otherwise there would be an open interval \( (x-t,x+t) \subseteq  [x, x + r)  \) for some \( t  > 0  \), which is impossible. We can see this by looking at the left-endpoint; that is, any open set around \( x  \) will not be contained within \( [x,x+r) \).
    \end{itemize}
\end{eg}

\begin{eg}[Computing Interiors of Sets in Different Topologies]
    Given \( A = (a,b), B = [a,b), C = [a,b],    \) and \( D= \Q  \). Find the interior of \( A,B,C,D \) as subsets of \( \R  \) endowed with 
   \begin{itemize}
       \item \textbf{The Euclidean Topology:}
               \( \mathring{A} = A  \), \( \mathring{B} = A    \), \( \mathring{C} = A  \), and \( \mathring{D} = \emptyset \);
           \item \textbf{The Lower Limit Topology:}
\( \mathring{A} = A   \), \( \mathring{B} = B  \), \( \mathring{C} = B  \), and \( \mathring{D} = \emptyset \);
\item  \textbf{The Cofinite Topology} \( \mathring{A} = \mathring{B} = \mathring{C} = \mathring{D} = \emptyset \)
   \end{itemize} 
\end{eg}

\subsection*{T4: Closed Sets and Closure}

\begin{eg}[Closed Sets in \( \R  \)]
        The sets \( [a,\infty), (-\infty ,a], [a,b]  \) for \( a,b \in \R  \) with \( a < b  \) are closed in \( \R  \) equipped with the Euclidean Topology. 

           It suffices to observe that 
           \begin{itemize}
               \item \( \R \setminus  [a,\infty) = (-\infty , a)  \) which is clearly open in \( \R  \) and so we see that \( [a,\infty) \) is closed in \( \R  \).
               \item \( \R \setminus  (-\infty , a] = (a,\infty )  \) which is clearly open in \( \R  \) and so we see that \( (-\infty , a]  \) is closed in \( \R  \).
               \item \( \R \setminus  [a,b] = (-\infty , a) \cup (b,\infty )  \) is a union of open sets in \( \R  \) and so the complement of \( [a,b] \) is open in \( \R  \); that is, \( [a,b] \) is a closed set in \(\R \).
           \end{itemize}
\end{eg}

   \begin{eg}[Non-examples]
       The sets \( [a,b) \), \( (a,b] \) for \( a,b \in \R  \) with \( a < b  \) are not closed in \( \R  \) equipped with the Euclidean Topology.
           \begin{itemize}
               \item Note that \( \R \setminus  [a,b) = (-\infty , a) \cup [b,\infty)  \) is not open because \( [b,\infty) \) is not an open set in \( \R  \). Indeed, we can find an open ball centered at \( b  \) such that it is not completely contained in \( [b,\infty) \). That is, \( b  \) is not an interior point of the complement. Hence, we conclude that \( [a,b) \) is not a closed set in \( \R  \).
               \item  Likewise, \(\R \setminus  (a,b] = (-\infty , a] \cup (b, \infty )  \) since \( a  \) is not an interior point of \( \R \setminus  (a,b] \).
            \end{itemize}
    \end{eg}

    \begin{eg}
       In a metric space, singletons are closed.
    \end{eg}

    \begin{eg}
        Which of the sets in the previous example are closed in \( \R  \) equipped with the lower limit topology?
        \begin{itemize}
            \item \( \R \setminus  [a,+\infty) = (-\infty ,a ) = \bigcup_{ n \in \N  }^{  } [a-n,a)   \) is open since we can write the complement of \( [a,\infty)  \) as a union of basis elements of \( {\mathcal{T}}_{L} \). That is, \( [a,+\infty) \) is closed in the lower limit topology.
            \item \( \R \setminus  (-\infty ,a] = (a,+\infty ) = \bigcup_{ n \in \N  }^{  } [a+ \frac{ 1 }{ n } , a + n) \) is open in \( \R  \) and so \( [a,+\infty) \).
            \item Similarly, \( \R \setminus  [a,b) = (-\infty , a) \cup [b,+\infty) \) is open as a union of open sets, so \( [a,b) \) is closed in the Lower Limit Topology.
            \item \( \R \setminus  (a,b] = (-\infty , a] \cup (b,\infty ) \) is not open since \( a  \) is not an interior point of the complement; that is, there does not exists an interval of the form \( [c,d)  \) with \( c,d \in \R  \) with \( c < d  \), \( a \in [c,d)  \) that is contained inn \( A  \). Thus, it follows that \( (a,b]  \) is not closed.
        \end{itemize}
    \end{eg}

\begin{eg}
    Which of the following sets are closed in \( \R^{2} \) equipped with the Euclidean topology? 
    \begin{enumerate}
        \item[(i)] \( C =  \{ ({x}_{1}, {x}_{2}) \in \R^{2} : {x}_{2} \leq 1  \}  \);

            Notice that \( \R^{2} \setminus  C  = \{ ({x}_{1}, {x}_{2}) \in \R^{2} : {x}_{2} > 1  \}  \). We can prove that this set is open in \( \R^{2} \) using the same methods found in homework 1. Hence, it follows that \( C  \) is closed in \( \R^{2} \).
        \item[(ii)] \( A = \{ 0 \}  \times \R  \).

            Again, we can show that \(  \R^{2} \setminus  A   \) is open.
    \end{enumerate}
\end{eg}

\begin{eg}[Dense Sets In A Topological Space]
   \begin{enumerate}
       \item[(i)] The set of rationals \( \Q  \) is dense in \( \R  \) equipped with the Euclidean Topology (and with the Lower Limit Topology)
        \item[(ii)] More in general, for \( n \in \N  \), the set \( \Q^{n} \) is dense in \( \R^{n} \) equipped with the Euclidean Topology.
        \item[(iii)] The set of irrationals \( \R \setminus  \Q  \) is dense in \( \R  \) with the Euclidean Topology.
   \end{enumerate} 
\end{eg}

\section*{T5: Limit and Boundary Points}

\begin{eg}
    In \( \R  \) equipped with the Euclidean topology find the derived of \( A = (-\infty , 1 ) \cup \{ 2  \}  \) and \( B = \{ n^{-1} : n \in \N  \}  \).
\end{eg}

\begin{eg}
    In \( \R  \) equipped with the Euclidean topology ind the boundary of the following sets
    \begin{enumerate}
        \item[(1)] \( A = [a,b]  \) where \( a,b \in \R  \) with \( a < b  \);
        \item[(2)] \( B = (a,b)  \) where \( a,b \in \R  \) with \( a < b  \);
        \item[(3)] \( C = \Q  \).
    \end{enumerate}
\end{eg}

\begin{eg}
   Let \( (X,d) \) be a metric space \( x \in X  \), \( r > 0  \). We have that  
   \[  \partial {B}_{r}(x) \subseteq \{ y \in X : d(y,x) = r  \}. \]
   Does equality necessarily hold?
\end{eg}

\begin{eg}[Finding the Interior, Closure, Derived, and Boundary]
   \begin{enumerate}
       \item[(1)] \( A =  \{ ({x}_{1}, {x}_{2}) \in \R^{2} : {x}_{2} \neq 0  \}  \)
        \item[(2)] \( B = \Big\{ ({x}_{1}, {x}_{2}) \in \R^{2} : \frac{ {x}_{1}^{2}  }{ a^{2}  }  + \frac{ {x}_{2}^{2} }{  b^{2} }  < 1 \Big\} \) where \( a,b \in \R  \).
        \item[(3)] \( C =  \{ ({x}_{1}, {x}_{2}) \in \R^{2} : {x}_{2} = \sin ({x}_{1}^{-1}) , {x}_{1} > 0  \}  \).
        \item[(4)] \( D = \Q \times \R  \)
   \end{enumerate} 
\end{eg}


\section*{Definitions}

\begin{definition}[Metric]
    Let \( X  \) be a nonempty set. A \textbf{metric} (or \textbf{distance}) on \( X  \) is a function \( d : X \times X \to [0,\infty)  \) such that for all \( x,y, z \in X  \)
    \begin{enumerate}
        \item[(i)] \( d(x,y) = 0  \) if and only if \( x = y  \)
        \item[(ii)] \( d(x,y) = d(y,x) \)
        \item[(iii)] \( d(x,y) \leq d(x,z) + d(z,y) \) 
    \end{enumerate}
    The pair \( (X,d) \) is called a \textbf{metric space}. 
\end{definition}

\begin{definition}[Open Ball in a Metric Space]
    Let \( (X,d) \) be a metric space, \( {x}_{0} \in X  \). We define the \textbf{open ball} centered \( {x}_{0} \) with radius \( r > 0  \) as 
    \[  {B}_{r}({x}_{0}) = \{ x \in X : d(x,{x}_{0}) < r  \}. \]
\end{definition}

\begin{definition}[Open Set in a Metric Space]
    Let \( (X,d) \) be a metric space. We say that a nonempty set \( A \subseteq X  \) is open if for all \( x \in A  \), there exists a \( \delta > 0  \) such that \( {B}_{\delta}(x) \subseteq A  \).
\end{definition}

\begin{definition}[Continuity In A Metric Space]
    Let \( (X,{d}_{X}) \) and \( (Y, {d}_{Y}) \) be two metric spaces. We say that a function \( f: X \to Y  \) is \textbf{continuous} at \( {x}_{0} \in X   \) if for all \( \epsilon > 0  \), there exists a \( \delta > 0  \) such that if \( {d}_{X}(x,{x}_{0}) < \delta \), then \( {d}_{Y}(f(x), f({x}_{0})) < \epsilon \). 

    We say that \( f  \) is continuous on \( X  \) if \( f  \) is continuous for all \( {x}_{0} \in X  \).
\end{definition}

\begin{definition}[Topology]
   Let \( X  \) be a nonempty set. A \textbf{topology} on \( X  \) is a collection of subsets of \( x  \) denoted by \( \mathcal{T} \) such that  
   \begin{enumerate}
       \item[(i)] \( \emptyset, X \in \mathcal{T} \);
        \item[(ii)] If \( {A}_{1}, {A}_{2} \in \mathcal{T} \), then \( {A}_{1} \cap {A}_{2} \in \mathcal{T} \);
        \item[(iii)] If \( \{ {A}_{i} \}_{i \in I} \), then \( \bigcup_{ i \in I  }^{  }  {A}_{i} \in \mathcal{T} \).
   \end{enumerate}
\end{definition}

\begin{definition}[Metrizable Topological Spaces]
    We say that a topological space \( (X,\mathcal{T}) \) is \textbf{metrizable} if there exists a metric \( d  \) that induces the topology \( \mathcal{T} \). 
\end{definition}


\begin{definition}[Basis of a Topology]
   Let \( (X,\mathcal{T}) \) be a topological space. A \textbf{basis} for \( \mathcal{T} \) is a collection of open sets \( \mathcal{B} \subseteq \mathcal{T} \) such that any set \( U \in \mathcal{T} \) can be written as a union of sets in \( \mathcal{B} \); that is, for all \( U \in \mathcal{T} \), 
   \[  U = \bigcup_{ B \in \mathcal{B} }^{  } B. \]
\end{definition}

\begin{definition}[Neighborhood]
   Let \( (X,\mathcal{T}) \) be a topological space and \( x \in X  \). A \textbf{Neighborhood} of \( x  \) is a subset \( N  \) of \( X  \) such that there exists \( U \in \mathcal{T} \) such that \( x \in U \subseteq N  \). The family of all neighborhoods of \( x  \) is called the \textbf{neighborhood system} of \( x  \) and is denoted by \( \mathcal{N}(x) \).
\end{definition}

\begin{definition}[Interior]
   Let \( X  \) be a topological space, \( A \subseteq X  \), and \( x \in X  \). We say that \( x  \) in an \textbf{interior point} of \( A  \) if there exists an open set \( U \subseteq X   \) containing \( x  \) such that \( x \in U \subseteq  \). 

   The \textbf{interior} of \( A  \) is the set of all interior points of \( A  \) denoted by
   \[  \text{int}(A) = \{ x \in A : \exists U \subseteq X \ \text{such that} \ x \in U \subseteq A \}.   \]
\end{definition}

\end{document}

